\documentclass[10pt,a4paper]{amsart}
\usepackage[margin=19mm]{geometry}
\usepackage{amsmath,amssymb,mathtools}
\usepackage[scr=rsfs,cal=euler]{mathalfa}
\usepackage{xfrac}
\usepackage[unicode,hidelinks,bookmarks=false]{hyperref}
\newcommand{\ie}{i.e.\ }
\newcommand{\cf}{cf.\ }
\newcommand{\resp}{respectively\ }
\newcommand{\Subsection}{\subsection}
\providecommand{\nobreakdash}{\nobreak\hbox{-}}
\setlength{\emergencystretch}{2em}
\multlinegap0pt
\usepackage{bm}
\usepackage{bbm}
\usepackage{dsfont}
\usepackage{xspace}
\usepackage{cancel}
\usepackage{mathtools}
\usepackage{amsthm}
\makeatletter
\@ifundefined{theorem}{\newtheorem{theorem}{Theorem}}{}
\@ifundefined{lemma}{\newtheorem{lemma}[theorem]{Lemma}}{}
\@ifundefined{proposition}{\newtheorem{proposition}[theorem]{Proposition}}{}
\makeatother
\def\al{\alpha}
\newcommand{\A}{{\mathcal{A}}}
\newcommand{\B}{{\mathcal{B}}}
\newcommand{\C}{{\mathcal{C}}}
\newcommand{\R}{{\mathcal{R}}}
\newcommand{\bC}{\mathbb{C}}
\newcommand{\eps}{\varepsilon}
\newcommand{\qand}{\quad\text{and}\quad}
\renewcommand{\S}{{\mathcal{S}}}
\title[Morita equivalence for quantum graphs]{Morita equivalence for quantum graphs}
\author{Alexandros Chatzinikolaou, Gage Hoefer, Nikolaos Koutsonikos-Kouloumpis,, Ioannis Apollon Paraskevas}
\date{Source: arXiv:2604.18065v1 (CC BY 4.0)}
\begin{document}
\maketitle

\noindent\emph{Source and licence.} Morita equivalence for quantum graphs. Version arXiv:2604.18065v1 (CC BY 4.0), distributed under the Creative Commons Attribution 4.0 International licence, as recorded in the arXiv licence metadata for this exact version and shown on the abstract page https://arxiv.org/abs/2604.18065v1. Full rights notice: licenses/arxiv-2604.18065-CC-BY-4.0.txt.\par\smallskip

\noindent\textit{Editorial scope.} Counted unit: the Proposition whose source label is \texttt{prop:reduced\_qgraph} in the article (source0.tex lines 1707--1709, source label \texttt{prop:reduced\_qgraph}), delivered together with the article's own complete original proof of it (source0.tex lines 1711--1864, one \texttt{proof} environment of 154 lines). No mathematics is written, repaired or re-derived: statement and proof are the article's own text line for line. Required context retained uncounted: lem:rij\_structure (lines 1630--1641). Every definition, display and label of those lines is retained unchanged. Omitted from this package and not redistributed: 1--1629, 1642--1706, 1710--1710, 1865--2918. Nothing inside a retained excerpt is omitted or altered: every retained body is delivered exactly as the article's lines are written, so no omission receipt of any kind accompanies this package. Citations: the retained excerpts carry no citation command at all, so no bibliography entry of the article is transcribed and no reference is redistributed. Reference adaptation: none was needed and none was made; every reference inside the delivered unit resolves inside it, so no unresolved reference or citation is produced. Printed slips kept literal: no printed word, symbol, hypothesis or number of the counted statement or of its proof was altered. Layout and class: the article's own class is \texttt{amsart} with packages including amssymb, amstext, amscd, amsmath, mathtools, paralist, color, dsfont, rotating, verbatim, enumerate, enumitem, float, setspace; the wrapper is the lane's pinned \texttt{amsart} editorial wrapper at 10pt on A4, which restates the theorem-like environments the retained text needs and adds the article's own macro definitions that the retained text needs (\texttt{\textbackslash A}, \texttt{\textbackslash B}, \texttt{\textbackslash C}, \texttt{\textbackslash R}, \texttt{\textbackslash S}, \texttt{\textbackslash al}, \texttt{\textbackslash bC}, \texttt{\textbackslash eps}, \texttt{\textbackslash qand}), with extra wrapper packages (bm, bbm, dsfont, xspace, cancel, mathtools). The delivered numbering is the unit's own. Assets: no retained span contains a figure environment, a graphics-inclusion command, a PDF-page inclusion, a table or a TikZ picture; no image file is copied into this package, the package declares no asset dependency and no image file is redistributed with it. Licence: version arXiv:2604.18065v1 of this article is distributed under the Creative Commons Attribution 4.0 International licence, as recorded in the arXiv licence metadata for this exact version and shown on the abstract page https://arxiv.org/abs/2604.18065v1. A full rights notice is delivered at licenses/arxiv-2604.18065-CC-BY-4.0.txt. Characters: every retained excerpt is pure ASCII, so the delivered engine, XeLaTeX with its default Latin Modern text font, reports no missing character.
\begin{lemma}\label{lem:rij_structure}
For each $i,j\in [k]$, let $\R_{ij}$ be defined via the slice construction described above. Then the following hold:
\begin{enumerate}
\item The space $\R_{ij}$ is independent of the choice of rank one operator used in its definition.
\item We have the decomposition
\[
\S_{ij}
=
M_{\alpha_i,\alpha_j}\otimes \R_{ij}.
\]
\end{enumerate}
\end{lemma}

\begin{proposition}\label{prop:reduced_qgraph}
The skeleton $(\R,\C)$ of an irreducibly acting quantum graph $(\S,\A)$ is a quantum graph and $(\S,\A)$ is a full pullback of $(\R,\C)$.
\end{proposition}

\begin{proof}
By Lemma~\ref{lem:rij_structure}, we have
\[
\S_{ij}=M_{\alpha_i,\alpha_j}\otimes \R_{ij}
\qquad \text{for all } i,j\in [k].
\]
We first show that $(\R,\C)$ is a quantum graph. 
We have that $\R$ is a selfadjoint subspace since $(\S_{ij})^*=\S_{ji}$, and it
contains $I_L$ as 
\[
L_{\eps_{11}^{(i,i)}}(I_{\alpha_i}\otimes I_{n_i})=I_{n_i}.
\]
Moreover, $\R$ is trivially a bimodule over $\C'=\bigoplus_i\mathbb C I_{n_i}$ and therefore
$(\R,\C)$ is a quantum graph. We now construct the pullbacks. Define $\theta \colon \C \to \A_{\S}'$ by
\[
\theta (\oplus_{i=1}^k x_i)
=
\oplus_{i=1}^k (I_{\alpha_i}\otimes x_i),
\]
and since $\S$ is an $\A'$-bimodule, by the definition of $\A_\S$ we obtain $ \A'\subseteq \A_{\S}$ and hence $\theta(\C)\subseteq \A$.
This is a unital $*$-homomorphism since it is the direct sum of the canonical
amplifications $x_i\mapsto I_{\alpha_i}\otimes x_i$ on each block. 

We will now obtain a Kraus representation for $\theta$. For each $i\in [k]$ and $r\in[\al_i]$ define the operators $ v_{i,r}: H \to L$ by 
\[
v_{i,r}:\begin{pmatrix}
\zeta_1\\
\vdots\\
\zeta_i\\
\vdots\\
\zeta_k
\end{pmatrix} \mapsto \begin{pmatrix}
0 \\
\vdots \\
(e_r^* \otimes I_{n_i})(\zeta_i)\\
\vdots\\
0
\end{pmatrix}
\]
with zeros in all but the $i$-th coordinate and $e_r^*(\xi)=\langle e_r,\xi \rangle$ for $ \xi \in \bC^{\alpha_i}$.   Thus  $v_{i,r} \colon H\to L$ vanishes off $H_i$ and  
\[
v_{i,r}(\sum_s e_s\otimes\xi_s)=\xi_r \quad  \text{for} \quad \sum_s e_s\otimes\xi_s\in H_i.
\]
In addition, we have
\[
v_{i,r}^* : \begin{pmatrix}
\xi_1\\
\vdots\\
\xi_i\\
\vdots\\
\xi_k
\end{pmatrix} \mapsto \begin{pmatrix}
0 \\
\vdots\\
e_r \otimes \xi_i\\
\vdots\\
0
\end{pmatrix}.
\]
Note that
\[
\sum_{r=1}^{\alpha_i} v_{i,r}^*v_{i,r}=p_i
\qand
\sum_{r=1}^{\alpha_i} v_{i,r}v_{i,r}^*=\alpha_i I_{L_i},
\]
and therefore
\[
\sum_{i,r} v_{i,r}^*v_{i,r}=I_H
\qand
\sum_{i,r} v_{i,r}v_{i,r}^*=\bigoplus_{i=1}^k \alpha_i I_{L_i}.
\]
Fix $x=\oplus_i x_i\in\C\subseteq \B(L)$.
For 
\[
\zeta= (\zeta_j)_{j=1}^k\in H \text{ with }  \zeta_j=\sum_{s}e_s\otimes\xi_s^{(j)}\in H_i,
\]
we have
\[
(v_{i,r}^*xv_{i,r})(\zeta)= v_{i,r}^* 
x_i \xi_r^{(i)} = 
e_r \otimes x_i \xi_r^{(i)},
\]
and hence
\[
\sum_{r=1}^{\alpha_i} (v_{i,r}^*xv_{i,r})(\zeta)
=
\sum_{r=1}^{\alpha_i} e_r\otimes x_i\xi_r^{(i)}
=
(I_{\alpha_i}\otimes x_i)\zeta.
\]
Summing over all $i\in [k]$ yields the Kraus form
\[
\theta_{}(x)=\sum_{i=1}^k\sum_{r=1}^{\alpha_i} v_{i,r}^*\,x\,v_{i,r}.
\]

To show the pullback relations, fix $i,j \in [k]$ and  let $y$ be $\R_{ij}$. We claim that
\[
v_{i,r}^*\,y\,v_{j,r'}=(\eps^{(i,j)}_{rr'}\otimes y),
\]
under the usual identification of $\B(H_j,H_i) \subseteq \B(H)$.
To verify this, it suffices to evaluate both sides on elementary tensors $e_s\otimes\eta\in H_j$
where $1\le s\le \alpha_j$ and $\eta\in L_j$. Indeed, we have
\[
v_{i,r}^*\,y\,v_{j,r'}(e_s\otimes\eta)
=
v_{i,r}^*(\delta_{r's}\,y\eta)
=
\delta_{r's}\,(e_r\otimes y\eta),
\]
and
\[
(\eps^{(i,j)}_{rr'}\otimes y)(e_s\otimes\eta)
=
(\eps^{(i,j)}_{rr'}e_s)\otimes y\eta
=
\delta_{r's}\,e_r\otimes y\eta.
\]
Since $y\in\R_{ij}$, we obtain
\[
v_{i,r}^*\,y\,v_{j,r'}=\eps^{(i,j)}_{rr'}\otimes y\in
M_{\alpha_i,\alpha_j}\otimes \R_{ij}=\S_{ij},
\]
showing that 
\[
v_{i,r}^*\R_{ij} v_{j,r'} \subseteq \S_{ij}.
\]
Conversely, if $x\in\S_{ij}$ write
\[
x=\sum_{\mu=1}^{\alpha_i}\sum_{\nu=1}^{\alpha_j} \eps^{(i,j)}_{\mu \nu}\otimes y_{\mu \nu}
\text{ for } y_{\mu\nu}\in\R_{ij}.
\]
Then for $\eta\in L_j$ we obtain
\[
v_{i,r}x v_{j,r'}^*(\eta)
=
v_{i,r}(x(e_{r'}\otimes\eta))
=
v_{i,r}(\sum_{\mu=1}^{\alpha_i} e_\mu\otimes y_{\mu r'}\eta)
=
y_{rr'}\eta,
\]
so
\[
v_{i,r} x v_{j,r'}^*=y_{rr'}\in\R_{ij}.
\]
Therefore, we have
\[
v_{i,r}^*\,\R\,v_{j,r'}\subseteq \S 
\qand
v_{i,r}\,\S\,v_{j,r'}^*\subseteq \R,
\]
and $(\S,\A)$ is a  pullback of $(\R,\C)$.
Finally $\theta$ is clearly injective and hence it is a full pullback.
\end{proof}

\end{document}
